\chapter{Subjective Probability and Its Dynamics}
\label{sec:p15-rationality-4-1}

Confidence in a lottery loss, tomorrow's cloud cover, and a lottery win can
differ even when none is treated simply as known or unknown. Alan H\'ajek
and Julia Staffel investigate the Bayesian proposal that these degrees of
belief, or \emph{credences}, are subject to probability-based norms. Their
chapter~4.1 of \emph{The Handbook of Rationality} distinguishes constraints
at a time from constraints on change through time, surveys their
justifications, and asks how the idealized account bears on imperfect
thinkers. The treatment here follows that complete 2021 presentation,
including its nine notes, summary table, and references
\parencite[pp.~239--251]{HajekStaffel2021Dynamics}.

This is a philosophical survey with mathematical components, not an
experiment establishing that people conform to its proposed norms.
An equation can establish a conditional formal result without establishing
its normative authority or its psychological applicability. The existing
finite-bet, conditioning, and prediction chapters supply exact neighboring
models; they are compared below rather than relabeled as this source's
complete theory.

\section{A calculus and its subjective interpretation}
\label{sec:p15-rat41-calculus}

\begin{hbdefinition}{def:p15-rat41-probabilism}{Finite-additive probability and probabilism}
Let \(\Omega\) be a nonempty space of possibilities and \(\mathcal F\)
an algebra of its events. A real-valued \(P\) on \(\mathcal F\) is a
\emph{finitely additive probability} when it is nonnegative, has
\(P(\Omega)=1\), and satisfies
\[
 P(X\cup Y)=P(X)+P(Y)\qquad(X\cap Y=\varnothing).
\]
For \(P(Y)>0\), its ratio conditional is
\[
                  P(X\mid Y)=\frac{P(X\cap Y)}{P(Y)}.
\]
\emph{Probabilism} is the additional normative thesis that a perfectly
rational agent's credences at a time are representable by such a
probability function. The calculus and this thesis are not identical.
Below \(C\) denotes a credence function; \(C_{\rm old}\) and
\(C_{\rm new}\) distinguish stages of an update.
\end{hbdefinition}

The algebra assumption makes the displayed operations available. In a
sentential formulation, normalization assigns one to a tautology and
additivity applies to logically incompatible sentences. Interpreting
sentences as events requires the truth conditions and equivalence policy
of that formulation; their syntax is not literally a set union.

The source explicitly declines to assume countable additivity. Its note~3
gives that stronger axiom and observes that classical syntax has no
arbitrary countable disjunction. On a finite space, the existing
Definition~\ref{def:p15-finite-conditioning} and
Theorem~\ref{thm:p15-finite-conditioning-laws} give the ratio construction
and successive-conditioning identity. The measure-theoretic development
in Part~\ref{part:p15} makes further assumptions that cannot be silently
imported into this finite-additive survey.

The same mathematical \(P\) can be interpreted as objective chance or as
subjective confidence. Frequencies, propensities, and best-system accounts
of chance are possibilities the authors mention, not equivalent definitions.
For a die assigned equal probabilities to six outcomes, the probability
of an even result is \(3/6=1/2\), and that of six conditional on an even
result is \((1/6)/(1/2)=1/3\). The arithmetic does not establish that a
physical die is fair or that an agent ought to use this model.

The authors' introductory range of confidence, from something existing
through lottery loss and cloud cover to lottery win and nothing existing,
is qualitative. Their note~1 stresses that passage from comparative
confidence to numbers itself needs an account, naming Fishburn and
Stef\'ansson. Notes~2 and~4 respectively identify conditional-probability
primitive approaches of Popper and R\'enyi and distinguish conditional
probability from the probability of a conditional. Neither is reduced to
the ratio formula here
\parencite[pp.~239--240 and p.~248, nn.~1--4]{HajekStaffel2021Dynamics}.

\section{Three arguments for probabilism and their additional premises}
\label{sec:p15-rat41-probabilism-arguments}

\subsection{Betting prices and Dutch books}

On the betting account, an agent's credence in an event equals the price
at which the agent would either buy or sell a contract paying one unit
if that event occurs and zero otherwise. The source's objections concern
the proposed identification, not the probability calculus. There may be
no unique buying-and-selling price. A bet on an event controlled by the
agent can change the event's prospects: a person unlikely to imitate a
chicken may do so for a sufficiently large payment. Changing the stakes
can change behavior without changing the original credence. An offer of
a high-stakes bet can also suggest that the other party has inside
information. Finally, utility for money, enjoyment of gambling, and
public bravado can affect a price jointly with confidence. Elicitation
can therefore alter what it purports to measure
\parencite[p.~240]{HajekStaffel2021Dynamics}.

A Dutch book is a portfolio of bets with net loss in every possible
outcome. The argument for probabilism has distinct steps: nonprobabilistic
normalized prices permit a sure-loss portfolio; susceptibility is a
rational defect; and prices are identified with credences. The last two
steps require philosophical justification. The converse matters too:
probability prices exclude such a portfolio under the specified contract
model. The exact finite equivalence, including reversible trades and no
transaction charges, is proved in
Theorem~\ref{thm:p17-finite-price-coherence}. It does not say that anyone
actually accepts every permitted portfolio.

For assessments on conditional events, with a bet called off when its
condition fails, the corresponding coherence condition, its agreement with
this finite theorem on unconditional events, and the objections above
applied to its two-sided stakes are in \cref{sec:p15-w02b-p15-038-dutch-book}.

\begin{hbexample}{def:p15-rat41-election-book}{The source's ten-cent election book}
Treat Democratic and Republican victories as incompatible events. The
source agent quotes \(1/2\) for each, but \(9/10\) for their union.
The agent buys the two separate unit-paying contracts for a total of
one dollar and sells the union contract for ninety cents. The net
initial payment is ten cents. In either victory the contracts held and
owed pay the same; if neither party wins, all pay zero. Thus the final
net loss is ten cents in every outcome.

The incompatibility premise is essential: without it the two separate
contracts need not equal the union contract. The example assumes the
same settlement conditions and unrestricted buying and selling at those
quotes. It is a hypothetical price schedule, not an election forecast.
\end{hbexample}

Critics distinguish a practical vulnerability from an epistemic defect
and observe that an agent can refuse the proposed package. The authors
report Skyrms's reply that the betting story dramatizes inconsistent
evaluations rather than predicting actual exploitation. Ramsey and de
Finetti supply the historical theorem attribution, with Kemeny for the
converse. The source directs readers to Skyrms, Vineberg, and H\'ajek
for fuller disputes. These are separate arguments, not further premises
proved by the ten-cent calculation
\parencite[pp.~240--241]{HajekStaffel2021Dynamics}.

\subsection{Representation is not automatically psychological identification}

Ramsey's approach uses preferences over gambles to calibrate utility and
then determine credences through utility comparisons. H\'ajek and Staffel
distinguish consistency requirements such as transitivity from structural
assumptions introduced to make the representation possible. They describe
Savage's theorem as jointly representing preferences through a probability
function and a utility class, with probability unique and utilities unique
up to positive affine transformations. They contrast Jeffrey's alternative
axioms and Bolker's representation, which need not determine a unique
probability function
\parencite[p.~241]{HajekStaffel2021Dynamics}.

An argument for probabilism still has to defend the preference axioms
as rational requirements and identify the representing weights with the
agent's credences. Representability alone does not settle that
identification: the source mentions alternative nonprobabilistic
representations and implausibly reinterpreted maximization stories.
An antirealist may treat numerical credences as artifacts of a
representation; a realist may regard them as genuine, even primitive,
mental states. The authors place Stef\'ansson and Eriksson--H\'ajek on
different sides of these issues rather than pretending that one theorem
settles them.

Theorem~\ref{thm:p17-finite-bet-representation} gives the existing finite
linear betting model, with its own assumptions. Theorem~\ref{thm:p17-finite-vnm}
uses supplied lottery probabilities. Neither is Savage's full
preference-derived construction, and neither proves that every
representing number is an actual credence. The works on which the
surveyed axiomatizations and the chapter's other reported theorems depend,
with the status of each reported claim in this handbook, are listed in
\cref{tab:p15-w02b-p15-022-dependent-sources}.

\subsection{Accuracy and its measure}

If credences are mental states answerable to truth, they can be evaluated
by their proximity to the truth rather than by financial outcomes.
The source describes the ideal as full confidence in truths and zero
confidence in falsehoods. Intermediate confidence is evaluated through
a chosen accuracy measure.

\begin{hbdefinition}{def:p15-rat41-brier}{Finite Brier inaccuracy and representable credence vectors}
Fix finitely many events \(A_1,\ldots,A_m\) on a finite nonempty world
space. At world \(w\), let \(v(w)\) be their truth vector, with
coordinate one when \(A_j\) occurs and zero otherwise. A credence vector
\(c\in[0,1]^m\) has \emph{Brier inaccuracy}
\[
            S(c,w)=\sum_{j=1}^m(c_j-v_j(w))^2.
\]
Its coherent representation set \(K\) is the convex hull of the vectors
\(v(w)\): the set of their averages under probability masses on the
worlds. Thus \(c\in K\) exactly when a probability on these worlds
has the stated event credences. This definition also handles overlapping
events; coherence is not merely summing their coordinates to one.
\end{hbdefinition}

\begin{hbtheorem}{thm:p15-rat41-brier}{proved=annex:p15.rat41.brier}{Strict propriety, accuracy dominance, and a planned update}
For one event assigned probability \(q\in[0,1]\), the expected Brier inaccuracy
of reporting \(x\in[0,1]\) is
\[
        q(1-x)^2+(1-q)x^2=q(1-q)+(x-q)^2.
\]
It is uniquely minimized at \(x=q\), including \(q=0,1\).
Absolute-error scoring instead has expected loss \(q+(1-2q)x\):
it favors an endpoint if \(q\ne1/2\), and every report ties if \(q=1/2\).

For the finite vector model, every \(c\notin K\) has a coherent
replacement strictly more Brier-accurate in every world. No \(c\in K\)
is strictly dominated in every world by any credence vector.
For a fixed prior and a finite partition of possible observations,
planning to report the old conditional event credences on every
positive-probability cell uniquely minimizes prior expected total Brier
inaccuracy among all planned reports on those cells.
\end{hbtheorem}

The source describes extensions beyond Brier scores, citing de Finetti,
Joyce, Leitgeb--Pettigrew, and Pettigrew. The theorem above fixes the
finite Brier case rather than asserting those extensions without their
conditions. Turning dominance into probabilism requires the further
normative premise that dominated credences are irrational and a defense
of the accuracy measure. Strict propriety alone does not select one
epistemically privileged score from all eligible scores. The source
identifies this problem and directs a different purported justification,
Cox's theorem, to another chapter in note~5
\parencite[pp.~241--242 and p.~248, n.~5]{HajekStaffel2021Dynamics}.

\section{Four further synchronic proposals}
\label{sec:p15-rat41-synchronic}

\subsection{Regularity and possibility}

Regularity proposes positive credence for every possible event. Logical,
metaphysical, and epistemic possibility produce different requirements.
The source's favorable arguments appeal to keeping open what the evidence
has not excluded, retaining the ability to learn, and avoiding a
\emph{semi-Dutch book}: a contract that can lose and cannot gain.
They do not establish that possibility and positive probability are
equivalent.

If a possible event \(X\) has credence zero, selling a unit-paying bet
on \(X\) for zero has net payoff \(-\mathbf1_X\): no gain is possible,
and one unit is lost if \(X\) occurs. Equivalently, buying the
complementary unit-paying bet for one gives that payoff. The zero price
for \(X\) makes these transactions fair under the proposed betting
identification. With such weak loss,
unlike a strict Dutch book, excluding exploitation is a stronger condition
than ordinary finite coherence. Zero-probability states remain permitted
in Theorem~\ref{thm:p17-finite-price-coherence}.

\begin{hbtheorem}{thm:p15-rat41-positive-atoms}{proved=annex:p15.rat41.positive-atoms}{An uncountable outcome space cannot give every singleton positive real mass}
For any nonnegative finitely additive normalized probability, the set
of measurable singleton outcomes with positive probability is at most
countable. Consequently, if uncountably many singleton outcomes are
measurable and regarded as possible, regularity for all those singletons
is impossible with real-valued probabilities. Countable additivity is
not needed for this conclusion.
\end{hbtheorem}

The source also raises a practical difficulty if infinite expected
utility is possible, as in Pascal-wager or St.~Petersburg discussions:
positive weight on that possibility can force an infinite valuation,
which fails to distinguish it from certainty. This depends on the
particular extended-utility model; it is not a finite-payoff theorem.
H\'ajek and Staffel cite H\'ajek's rebuttals and other opponents of
regularity. Note~6 directs the decision-theoretic paradoxes to a separate
treatment. No blanket regularity requirement is adopted here
\parencite[pp.~242--243 and p.~248, n.~6]{HajekStaffel2021Dynamics}.

\subsection{Chance: two formulations of the Principal Principle}

The original Lewis formulation concerns a reasonable \emph{initial}
credence function \(C\). If \(X\) says that the chance of \(A\) at
time \(t\) is \(x\), and \(E\) is compatible admissible information,
it proposes
\[
                         C(A\mid X\cap E)=x.
\]
Admissibility requires that the information's influence on the outcome
credence operate through information about chances. The ratio expression
requires \(C(X\cap E)>0\); another conditional-probability framework
would need its own domain conditions.

The source example assigns heads probability one half given a noon
chance of one half and a heads result at 11:59. The earlier result is
treated as admissible; advance information about the future outcome
would not be. The principle does not demand deference to chance when
one already knows the outcome.

Lewis's 1994 discussion instead speaks of rational credence with
evidence limited to the past and present, not only an initial
evidence-free function. The authors expressly regard these formulations
as nonequivalent. Note~7 explains a route from the original to later
credences by conditioning on admissible evidence, but conditioning
itself is disputed; it also identifies Lewis's distinct New Principle.
Pettigrew's accuracy and almost-Dutch-book arguments are surveyed
justifications, not premises that make these formulations identical
\parencite[p.~243 and p.~248, n.~7]{HajekStaffel2021Dynamics}.

\subsection{Reflection and the reliability of one's future self}

Reflection proposes deference to future credence:
\[
               C_t(A\mid C_{t+n}(A)=x)=x\qquad(n>0),
\]
whenever the stated conditional is defined. In the source's car example,
being certain now that tomorrow's engine rattle will make one highly
confident that servicing is needed calls, under Reflection, for that
confidence now. The point is anticipated confidence, not observation
of the rattle already made.

The authors survey van Fraassen's diachronic Dutch-book argument and
Easwaran's expected-accuracy argument. Their essential qualification is
that unrestricted deference is implausible when the future self is
expected to become irrational or forgetful. They name Arntzenius,
Christensen and Talbott for objections and van Fraassen's later
modifications. The principle needs a restricted admissible situation
class; its displayed formula does not supply one by itself
\parencite[p.~243]{HajekStaffel2021Dynamics}.

\subsection{Indifference, entropy, and representation dependence}

The proposed Indifference Principle assigns \(1/n\) to each of \(n\)
possibilities when no evidence discriminates among them. The source's
lottery entrant, among one hundred participants without information
about winner selection, is accordingly assigned \(1/100\). This needs
an exhaustive, mutually exclusive set of relevant alternatives and the
normative premise about evidential symmetry; ignorance alone does not
specify a privileged partition.

For a finite probability vector \(p\), define
\[
       H(p)=-\sum_i p_i\log_2p_i,
       \qquad I(p)=\sum_i p_i\log_2p_i=-H(p),
\]
with \(0\log_20=0\). The second is the source's information convention.

\begin{hbtheorem}{thm:p15-rat41-indifference}{proved=annex:p15.rat41.indifference}{Finite entropy and minimax Brier properties of uniformity}
On \(n\geq1\) exhaustive outcomes, the uniform probability vector
uniquely maximizes \(H\), with value \(\log_2n\), and minimizes \(I\).
For Brier inaccuracy scored on the \(n\) outcome indicators, it also
uniquely minimizes the worst-world score, whose minimum is \(1-1/n\).
These conclusions concern the fixed finite outcome description.
\end{hbtheorem}

The general maximum-entropy proposal adds constraints representing
background information and chooses an entropy maximizer subject to
them. The declared optimization problem, its existence and uniqueness
conditions, and a reconstruction of Jaynes's argument for it are in
\cref{sec:p15-w02b-p15-043-maxent}. H\'ajek and Staffel name Jaynes as
an advocate, but note that continuous representations can conflict. In their factory example,
cubes have side lengths up to one foot. Uniformity in side length
assigns probability \(1/2\) to length at most half a foot; uniformity
in volume assigns \((1/2)^3=1/8\); uniformity in face area assigns
\((1/2)^2=1/4\). These are different priors, not inconsistent calculations
within a single stipulated prior. The source regards the lack of a
privileged symmetry description as a disputed objection, pointing to
Bertrand, van Fraassen, Norton, White and Novack.
In entropy terms, each of the three uniform priors is the unconstrained
maximum-entropy law only relative to Lebesgue measure in its own
coordinate. \Cref{thm:p15-w02b-p15-042-power-coordinates} computes the
continuous entropy of each prior in each of the three coordinates, and the
comparison in \cref{sec:p15-w02b-p15-042-relationships} tabulates the
values; the computation makes the conflict exact and does not choose a
privileged coordinate.
The existing reparameterization example in
Theorem~\ref{thm:p15-bayes-transformation} likewise distinguishes
transporting a fixed prior from reimposing uniformity in new coordinates
\parencite[pp.~243--244]{HajekStaffel2021Dynamics}.
\Cref{thm:p15-w02b-p15-044-cubes}
computes the factory example for side length, face area and volume, with
and without a declared mean, and \cref{thm:p15-w02b-p15-044-reference}
gives the effect of such a change of reference measure on the maximizer.

\section{Learning with and without certainty}
\label{sec:p15-rat41-dynamics}

\subsection{Conditioning as a proposed norm, not merely a ratio identity}

If an agent learns \(E\) with certainty and learns nothing additional,
conditionalization proposes
\[
                     C_{\rm new}(A)=C_{\rm old}(A\mid E).
\]
This demands \(C_{\rm old}(E)>0\) in the ratio framework and retention
of the relevant conditional relationships. The source neighbor example
has old confidence \(0.9\) in being home conditional on lights being on;
learning that the lights are on calls for new confidence \(0.9\).
The other stated confidence, \(0.5\) without that evidence, does not
determine the prior chance of the lights being on. No missing joint
distribution is inferred from those two numbers.

For fixed events with positive intersection probability, successive
conditioning gives the same result as conditioning on the intersection,
in either order, by Theorem~\ref{thm:p15-finite-conditioning-laws}.
This does not license order invariance for changed observation mechanisms,
information loss, or null intersections. The existing posterior-kernel
and sequential results in
Theorems~\ref{thm:p15-bayes-posterior-kernel} and
\ref{thm:p15-bayes-sequential-update} state their own model and path
regularity assumptions. The predictive construction in
Theorem~\ref{thm:p15-bayes-predictive-kernel} additionally integrates a
specified future observation law against a posterior. An update norm
does not supply that future law or prove that it describes a physical
process.

The authors survey diachronic Dutch-book arguments by Lewis and Skyrms,
including a converse only when no other Dutch-bookable offenses occur.
They preserve Christensen's comparison with differences between spouses,
Levi's objection that an agent aware of the strategy can decline initial
bets, and Skyrms's response involving a fully specified decision problem.
Their differences concern normative interpretation and betting
conditions, not whether the ratio identity is an algebraic identity.

Expected-accuracy arguments evaluate an update plan from the old
credences. The Brier planning conclusion of
Theorem~\ref{thm:p15-rat41-brier} is an explicit finite instance. The
source's objection is that after learning, the old credences are known
to be outdated, while any coherent new credences endorse themselves
under a strictly proper score. An old evaluation's preference for a plan
does not by itself establish that later compliance is rationally required.
The authors recount Pettigrew's distinction between planning to
conditionalize and being required to follow through, and cite other
expected-accuracy treatments rather than claiming one uncontested proof
of the full normative rule
\parencite[p.~244]{HajekStaffel2021Dynamics}.

\subsection{Jeffrey conditioning: new weights and retained conditionals}

The source's poor-light tablecloth example motivates learning without
certainty: seeing the cloth can increase confidence that it is blue
without making that confidence one. Whether some other proposition is
learned with certainty in such a situation is itself disputed.

\begin{hbdefinition}{def:p15-rat41-jeffrey}{A finite Jeffrey update}
Let \(E_1,\ldots,E_n\) partition a finite world space, with
\(C_{\rm old}(E_i)>0\) for every cell. Let new cell weights
\(r_i\geq0\) sum to one. The Jeffrey update is
\[
 C_{\rm new}(A)=\sum_i r_i C_{\rm old}(A\mid E_i).
\]
The input weights \(r_i\) come from the learning situation; the formula
does not calculate them from an unspecified experience. The retained
within-cell conditional relationships are the further stability premise
often called rigidity. On a cell receiving zero new weight, no new
ratio conditional is defined by that formula.
\end{hbdefinition}

\begin{hbtheorem}{thm:p15-rat41-jeffrey}{proved=annex:p15.rat41.jeffrey}{Existence, rigidity, certainty, and noncommutativity}
The finite Jeffrey formula defines a probability with cell masses
\(r_i\). On every positive new-weight cell it retains the old
conditional law. It is the unique probability with those weights and
that retained law on positive new-weight cells. A point-mass choice of
weights reduces it to ordinary conditioning on the selected cell.

For a two-by-two table of worlds initially assigned probabilities
\[
             \begin{pmatrix}2/5&1/10\\1/10&2/5\end{pmatrix},
\]
let \(E\) be the first row and \(F\) the first column. Resetting
\(E\)'s probability to \(4/5\) and then \(F\)'s to \(4/5\) gives
\[
             \begin{pmatrix}64/85&1/10\\4/85&1/10\end{pmatrix}.
\]
Resetting \(F\) first and \(E\) second gives instead
\[
             \begin{pmatrix}64/85&4/85\\1/10&1/10\end{pmatrix}.
\]
Thus these two Jeffrey reset operations do not commute.
\end{hbtheorem}

The numeric table is an editorial realization of the source's
noncommutativity point. After the first sequence, \(F\) has its latest
weight \(4/5\), but \(E\) has changed again, to \(29/34\). One must
not interpret sequential resets as guaranteeing preservation of both
constraints. The authors describe diachronic Dutch-book support for
Jeffrey conditioning by Armendt and Skyrms and report that no
accuracy-based argument was known to them at the time of writing.
That report is not a theorem of impossibility or a current literature
census
\parencite[p.~245]{HajekStaffel2021Dynamics}.

\subsection{Minimal change and its orientation}

For probability vectors \(Q=(q_j)\) and \(P=(p_j)\), the source defines
relative entropy by
\[
                    D(Q,P)=\sum_{j:q_j>0}q_j\log_2(q_j/p_j),
\]
with value infinity if some \(q_j>0=p_j\). The zero contribution is
defined by continuity. Order matters: \(D(Q,P)\) is not generally
\(D(P,Q)\). In the following result the new distribution is the first
argument and the old one the second.

\begin{hbtheorem}{thm:p15-rat41-minimal-update}{proved=annex:p15.rat41.minimal-update}{Finite partition constraints and three measures of change}
Under the Jeffrey assumptions, write \(P=C_{\rm old}\) and let
\(Q^*\) be the Jeffrey distribution with target masses \(r_i\).
Among all probability masses \(Q\) with \(Q(E_i)=r_i\),
\(Q^*\) uniquely minimizes \(D(Q,P)\). It also uniquely minimizes
\(\sum_j(\sqrt{q_j}-\sqrt{p_j})^2\).

It minimizes total variation \(\frac12\sum_j|q_j-p_j|\), with minimum
\(\frac12\sum_i|r_i-P(E_i)|\), but that minimizer need not be unique.
For prior masses \((1/4,1/4,1/2)\) and the constraint assigning one
to the first two worlds together, both \((1/2,1/2,0)\) and
\((3/4,1/4,0)\) attain total variation \(1/2\).
\end{hbtheorem}

Certainty is the single-cell-weight case of this theorem. The result
shows how particular measures implement conservative change; it does
not establish why minimal departure from outdated credences should be
an epistemic duty or why one measure is privileged. It also makes
explicit a nonuniqueness that must not be hidden by calling a
variation-distance update \emph{the} closest distribution.

The source's general dynamics problem imposes an arbitrary constraint
on an old credence function and seeks a closest admissible new one.
Examples include conditional-credence requirements or an interval
constraint on an event. Existence and uniqueness must be established
for the constraint and measure in use; they do not follow from the
problem's wording. The authors point to van Fraassen's Judy Benjamin
problem as an objection for a conditional-credence constraint, without
giving its full construction in this chapter
\parencite[pp.~244--245]{HajekStaffel2021Dynamics}.

\section{Uniqueness and non-ideal thinkers}
\label{sec:p15-rat41-nonideal}

The surveyed norms do not in general determine a single rational
credence function. \emph{Permissivism} allows more than one rational
doxastic response to a body of evidence. \emph{Intrapersonal uniqueness}
allows only one response for a particular thinker with that evidence;
\emph{interpersonal uniqueness} makes the same response uniquely rational
for every thinker with it. Even acceptance of uniqueness does not supply
a method of finding that response. The source identifies competing
defenses and survey treatments, and note~8 locates Carnap and objective
Bayesianism in this debate
\parencite[pp.~245--246 and p.~248, n.~8]{HajekStaffel2021Dynamics}.

\subsection{Imprecision and information loss}

H\'ajek and Staffel question whether point-valued credences accurately
represent ordinary confidence or fit very unspecific evidence. Their
cloudy-Vladivostok example suggests a broad interval such as
\([0.1,0.9]\) rather than a uniquely determined number. A set of precise
probabilities can represent an imprecise state, but exact endpoints can
introduce a new artificial precision. Norms, their justifications, and
decision rules then need reformulation. The authors name Joyce and Levi
for motivations and Elga and Rinard for a decision-theoretic dispute.
The imprecise probability treatment in
Chapter~\ref{sec:p15-imprecise-probability} specifies its own lower-expectation
and conditioning framework; it is not merely replacing every scalar
here by an interval.

Conditioning-only models also retain certainty once acquired. Forgetting
and information loss need not be irrational, and expected forgetting
already limits Reflection. The source points to Titelbaum's framework
for information loss and self-location, including who, where, and when
the thinker is, and to Arntzenius's examples. Those are separate belief-change
models: the old conditionalization formula is not claimed to implement
them
\parencite[p.~246]{HajekStaffel2021Dynamics}.
No chapter of this handbook develops Titelbaum's framework or a model of
self-locating credence, so these routes have no handbook treatment to which
they can be mapped. The nearest finite construction is imperfect recall in
extensive-form decision problems
(\cref{def:p17-rat91-memory,thm:p17-rat91-memory}); it concerns plans and
payoffs under forgetting and absentmindedness and does not supply a credal
update rule for information loss or for self-location.

\subsection{Approximating coherence is measure-dependent}

Assigning every tautology credence one exceeds human reasoning
capacities. The source also cites Slovic--Lichtenstein's 1971 review for
reported departures in response to evidence and describes a dispute
about their interpretation. The selected chapter supplies no task,
sample, design, or quantitative result for that attribution. It is
therefore not a validated population estimate or a finding that people
are generally rational or irrational. Such a claim needs the primary
empirical record required by Method~\ref{def:p19-evidence-record}.

If ideal norms are treated as regulative, two questions remain: how to
measure approximation, and why a specified improvement is beneficial.
The source's incoherent pair \((0.3,0.4)\) for a proposition and its
negation illustrates this problem.

\begin{hbtheorem}{thm:p15-rat41-graded}{proved=annex:p15.rat41.graded}{A finite, explicitly scored improvement toward coherence}
For the truth-vector model of Definition~\ref{def:p15-rat41-brier},
let \(y\in K\) minimize squared Euclidean distance from \(c\notin K\),
and put \(c_s=(1-s)c+sy\), \(0<s\leq1\). Then \(c_s\) is strictly
more Brier-accurate than \(c\) in every world, and its distance to
\(K\) is \((1-s)\|c-y\|\).

If allowed betting-stake vectors have Euclidean norm at most one,
the largest guaranteed Dutch-book loss at quote vector \(c\) is
exactly \(\|c-y\|\). At \(c_s\) it is \((1-s)\|c-y\|\).
Thus the improvement has both benefits under these declared measures
and this stake normalization, not under every conceivable measure.

For the source pair \(c=(0.3,0.4)\), the closest coherent pair is
\(y=(0.45,0.55)\), with squared distance \(0.045\).
The absolute-distance minimizers are every pair \((a,1-a)\) with
\(0.3\leq a\leq0.6\). If instead one explicitly minimizes
\(\sum_i q_i\log_2(q_i/c_i)\) over probability pairs \(q\), the
unique minimizer is \((3/7,4/7)\approx(0.429,0.571)\). The same pair is
the unique minimizer of the reversed orientation
\(\sum_i c_i\log_2(c_i/q_i)\), and adding \(\sum_i(c_i-q_i)\) to the
first objective does not change it. These expressions are extensions to
an unnormalized positive reference vector, not ordinary relative
entropy between two probability distributions.
\end{hbtheorem}

The different repairs in the theorem preserve the source's point that
metric choice matters. A coherent complementary pair must sum to one,
whatever the measure used to select it. Moreover, a relative-entropy
repair of an unnormalized vector needs an explicitly stated extension
and orientation; its name alone does not determine an answer.
The source itself reports the Kullback--Leibler repair of
\((0.3,0.4)\) as \(C^*(p)=0.438\), \(C^*({\sim}p)=0.572\)
\parencite[p.~247]{HajekStaffel2021Dynamics}. As printed these values
sum to \(1.010\), so they are not a coherent complementary pair. They
also differ from \((3/7,4/7)\approx(0.429,0.571)\), the unique
minimizer of each orientation in \cref{thm:p15-rat41-graded}. The
chapter states neither the formula nor the orientation behind its
figures, so the difference cannot be traced from the source; the
values in the theorem are those that follow from the stated
objectives.
The authors cite De Bona--Staffel and earlier treatments
for broader results and stress the importance of matching distance,
accuracy score, and betting normalization. Their reported logarithmic-score
comparison is not converted into a general Brier improvement theorem
\parencite[pp.~246--247]{HajekStaffel2021Dynamics}.

\subsection{Disagreement and higher-order evidence}

Two people calculate different answers to an agreed equal division of
a lunch bill. At least one calculation is wrong even if rational
responses to uncertain evidence may sometimes vary. The source
contrasts conciliatory reduction of confidence with steadfastness
conditional on actually assessing the evidence correctly. The latter
condition is not proof that a disputant can recognize being correct.

Not every disagreement has the same evidential structure. The source
distinguishes a better-informed interlocutor, a less-informed one,
different permissible starting priors, and independently agreeing
high-confidence judgments. These possibilities can require different
responses and constrain permissible priors. They do not justify automatic
averaging of every disagreement or counting dependent agreement as new
evidence. Note~9 points to Matheson and the Feldman--Warfield and
Christensen--Lackey anthologies, preserving this as a further debate
\parencite[p.~247 and p.~248, n.~9]{HajekStaffel2021Dynamics}.

\section{The source's justification map and reference use}
\label{sec:p15-rat41-reference-use}

Table~\ref{tab:p15-rat41-justifications} reproduces the intellectual
mapping in the source's Table~4.1.1 through independently worded row
descriptions. An entry records a \emph{putative} justification known
to the authors at writing, not independent acceptance of the argument.
The source gives two readings of an empty cell: the missing
justification may be impossible, or it may be a historical accident that
nobody has yet proved the required theorem, or perhaps even tried; it does
not say which reading applies to which cell. It also leaves the table open,
with a final row and a final column of ellipses, and says the table will
grow as further norms and kinds of justification are proposed
\parencite[pp.~247--248]{HajekStaffel2021Dynamics}. The table is not a survey
of literature published after the source.

\begin{table}[htbp]
\centering\small
\begin{tabular}{lcccccc}
\hline
Proposed norm & Betting & Accuracy & Representation & Entropy & Distance & \(\ldots\)\\
\hline
Probabilism & yes & yes & yes & & & \\
Regularity & yes & & & & & \\
Principal Principle & yes & yes & & & & \\
Reflection & yes & yes & & & & \\
Indifference & & yes & & yes & & \\
Conditioning & yes & yes & & & yes & \\
Jeffrey conditioning & yes & & & & yes & \\
\(\ldots\) & & & & & & \\
\hline
\end{tabular}
\caption{The 2021 source's recorded routes of putative justification.}
\label{tab:p15-rat41-justifications}
\end{table}

\begin{hbmethod}{def:p15-rat41-update-selection}{Select a probability-dynamics claim by its inputs and standing}
State the event language and possibility space, the interpretation of
the numbers, and whether the question concerns ideal coherence,
evidence-sensitive change, a representation, or actual human behavior.
Give prior support and the learning conditions rather than selecting an
update from its name.

For certainty about a positive-prior event with no further learning,
ordinary conditioning is the relevant proposed norm and mathematical
construction. For new weights on a partition, specify both the weights
and retained within-cell conditionals before using a Jeffrey update.
For other constraints, specify the objective's orientation and establish
feasibility and uniqueness rather than assuming every minimal-change
problem has one answer. At a prior-null event the ratio is unavailable;
at a null-valued observation a regular conditional version requires
the separate kernel and observation model.

For an assertion of rationality, identify the extra normative bridge:
betting vulnerability, accuracy dominance, preference axioms, deference,
symmetry, or conservative change. Keep the source's objections attached
to that bridge. For approximate repairs, declare the score, distance,
and stake normalization. For psychological claims, supply the actual
task and evidence record; numerical elegance supplies neither.
\end{hbmethod}

The source ends by identifying continuing questions, not by settling
them: quantitative versus comparative confidence; credences versus
binary belief and knowledge; imprecision, disagreement, and human
limitations; rival decision theories; and alternative representations
such as Dempster--Shafer, ranking, belief-revision, and possibility
frameworks. These are not interchangeable versions of ratio conditioning,
and the handbook treats each with its own representation and change rule.
Dempster--Shafer belief and plausibility functions are defined in
\cref{def:p15-r47-mass} (\cref{sec:p15-r47-evidence}), and conditioning a
lower prevision on an observed event is \cref{sec:p15-ip-conditioning}.
Ranking functions and their revision are \cref{sec:p15-rationality-5-3};
the finite bridge from negative ranks to possibility and necessity measures
is \cref{def:p15-r47-ranking,thm:p15-r47-ranking}, and possibilistic belief
statements are \cref{sec:p15-r47-modal}. AGM revision of a deductively
closed belief set is \cref{def:p08-agm}, with its faithful-preorder
representation in \cref{thm:p08-agm-representation}; ranked models and
rational closure for defeasible conditionals are
\cref{sec:p08-ranked-models-rational-closure}. None of these change rules
is the ratio formula applied to different numbers.
The present treatment preserves that open comparison boundary. Its finite
proofs establish their displayed claims; the competing norms retain
their philosophical status
\parencite[pp.~247--251]{HajekStaffel2021Dynamics}.

\index{subjective probability}\index{probabilism}\index{conditionalization}
\index{Jeffrey conditioning}\index{probability!regularity}
\index{Principal Principle}\index{Reflection Principle}
\index{Brier score}\index{credence!approximate coherence}
